\begin{abstract}
Data curation is universal in LLM training pipelines, yet practitioners choose filtering thresholds heuristically---leading to either wasted compute on noisy data or lost performance from over-filtering valuable diversity. We present the first controlled dose-response study of curation parameters for language model pretraining, treating perplexity thresholds as continuous variables and mapping their effect surfaces under fixed-token experimental design. Our key insight is that perplexity filtering exhibits a concave relationship with benchmark performance: intermediate thresholds near p50 outperform both extremes because they balance noise removal against diversity preservation. Experiments on GPT-2 variants trained on RedPajama-v2 confirm non-monotonic dose-response (quadratic $R^2=0.985$, interior peak at p44.5), validate the underlying mechanism through convergence dynamics analysis, and demonstrate 1.32\% improvement over industry-standard defaults at proof-of-concept scale. Scale transfer analysis shows optimal thresholds at 125M transfer to 1B models with ratio 0.85, enabling efficient small-scale optimization with predictable extrapolation. Our methodology transforms data curation from heuristic choice to principled optimization.
\end{abstract}
